\newif\ifglobal

%%%%%%%%%%%%%%%%%%%%%%%
%%% Compile to view %%%
%%%%%%%%%%%%%%%%%%%%%%%

%%%%%%%%%%%%%%%%%%%%%%%%%%%%%%%%%%%%%%%%%%%%%%%%%%%
%%% Readme:                                     %%%
%%% https://www.overleaf.com/read/bwdjvfjksvjy/ %%%
%%%%%%%%%%%%%%%%%%%%%%%%%%%%%%%%%%%%%%%%%%%%%%%%%%%

% >>> M?: marks a module setting number or important notice

% >>> M4: Set {./relative-path} to external files for this module <<<
\def \modulefiles {./31-RSADS}
% To add an external file, use:
% (Replace '---' with actual filename)
% '\input{\modulefiles/tables/---.tex}' for tables
% '\includegraphics{\modulefiles/figures/---}' for figures
% '\subfile{\modulefiles/---}' for register files


% >>> M1: This module is in global project? <<<
% 'YES' - uncomment; 'NO' - comment out
\globaltrue

\ifglobal
    \documentclass[main\_\_CN.tex]{subfiles}

\else
    % Font "FandolSong-Regular" used by ctex does not have GSUB table.
    % It causes warnings that do not affect compilation and can be ignored.
    \PassOptionsToPackage{quiet}{fontspec}

    \documentclass[openany, 10pt]{book}

    % >>> M2: In ./00-shared/preamble-trm-module, also find
    %         settings for visibility of todo notes, watermark <<<
    \usepackage{preamble-trm-module}

    % Enable Chinese typesetting
    \usepackage{ctex}

    % >>> M3: Temporary labels in Chinese
    %         you can add your own labels there <<<
    \myexternaldocument{./00-shared/temp-labels-cn}

\fi %\ifglobal


\setcounter{secnumdepth}{4}


% >>> M5: If this module needs any updates in
% './00-shared/preamble-trm-module.sty'
% (include additional package, etc.), contact your module PM
% or create the file './00-shared/changelog.tex' make the
% necessary updates yourself and document them in this file <<<

% Make text left-aligned and allow latex to break pages naturally, comment out for full justification
\raggedright
\raggedbottom


\begin{document}


% Sets language into which pre-defined bits of text are translated
\selectlanguage{chinese}
\setchineseformatting


% Do not delete this block
\ifglobal
% Add the following front matter if
% this module is outside of global project
\else
    \InputIfFileExists{readme}{}{}
    \clearpage
    \listoftodos
    \listofdonetodos
    \tableofcontents
    %\listoftables
    %\listoffigures
\fi


%%%%%%%%%%%%%%%%%%%%%%%%%
%%% TRM Module Begins %%%
%%%%%%%%%%%%%%%%%%%%%%%%%

\chapter{RSA 数字签名外设 (RSA\_DS)}
\label{mod:rsads}
\hypertarget{digsig}{}\hypertarget{rsads}{}

\section{概述}

RSA 数字签名外设 (RSA\_DS) 使用密码学算法，用于验证消息的真实性和完整性。该算法也可用于比如向服务器验证设备身份。

\chipname{} 包含一个 RSA 数字签名外设 (RSA\_DS)，可提供硬件加速，完成基于 RSA 的数字签名操作。HMAC 外设作为密钥导出函数 (Key Derivation Function, KDF)，使用存储在 eFuse 或密钥管理器生成的密钥作为输入密钥，输出 \begin{math}RSA\_DS\_KEY\end{math}。随后，RSA\_DS 外设使用 \begin{math}RSA\_DS\_KEY\end{math} 解密预先加密的参数，计算出签名。上述过程都发生在硬件层面，因此，在计算过程中，不论是解密 RSA 参数的密钥，还是用于 HMAC 密钥导出函数的
输入密钥，都对用户不可见。

\section{主要特性}

\chipname{} 的 RSA\_DS 外设具有以下特性：

\begin{itemize}
\item 硬件实现 RSA 签名运算 \begin{math}Z = X^{Y} \bmod M\end{math}，私钥相关运算子最大 4096 位
\item 私钥参数以密文存放在 flash 中，在 RSA\_DS 外设内以 \begin{math}RSA\_DS\_KEY\end{math} 解密，软件不可见明文
\item 对私钥数据及相关参数使用 SHA-256 计算摘要，用于防篡改检测
\item \begin{math}RSA\_DS\_KEY\end{math} 由硬件提供（对用户不可读），来源为基于 eFuse 中 \begin{math}HMAC\_KEY\end{math} 的 HMAC-SHA256 导出，或由密钥管理器部署
\end{itemize}

\section{功能描述}
\subsection{概述}

RSA\_DS 外设计算 RSA 签名的方式为 \begin{math}Z = X^{Y} \bmod M\end{math}，其中 \begin{math}Z\end{math} 是签名，\begin{math}X\end{math} 是输入消息，\begin{math}Y\end{math} 和 \begin{math}M\end{math} 是 RSA 私钥参数。

私钥参数以密文形式存储在 flash 中。用于解密私钥参数的密钥 (\begin{math} RSA\_DS\_KEY \end{math}) 可以来自于 HMAC 模块或密钥管理器模块。当该密钥：
\begin{itemize}
    \item 由 RSA\_DS 外设通过 HMAC 模块计算获取时，生成密钥所需的 HMAC 输入 \begin{math} HMAC\_KEY \end{math} 存放在 eFuse 中，只允许被 HMAC 模块访问。也就是说，只有 RSA\_DS 硬件才能解密私钥密文，软件无法获取私钥明文。关于 eFuse 和 HMAC 的相关细节请参照章节 \ref{mod:efuse} \textit{\nameref{mod:efuse}} 和章节 \ref{mod:hmac} \textit{\nameref{mod:hmac}}。
    \item 由 RSA\_DS 外设通过密钥管理器模块获取时，该密钥必须以安全的方式由密钥管理器部署而来。
\end{itemize}
需要签名时，软件直接将输入消息 \begin{math}X\end{math} 发送到 RSA\_DS 外设。RSA 签名计算完成后，软件将读取签名结果 \begin{math}Z\end{math}。

为方便描述，此处约定几个符号，它们的作用域仅限于本章。

\begin{itemize}

    \item {\bfseries 1\begin{math}^s\end{math}}:
        表示一个完全由 “1” 组成的长度为 \begin{math}s\end{math} 位的位串。

    \item {\bfseries [\begin{math}x\end{math}]\begin{math}_{s}\end{math}}:
       一个长度为 \begin{math}s\end{math} 位的位串，要求 \begin{math}s\end{math} 为 8 的整数倍。如果 \begin{math}x\end{math} 是一个数 (\begin{math}x\end{math} < 2\begin{math}^s\end{math})，那么其在位串中遵循小端字节序。\begin{math}x\end{math} 可以是一个变量，例如 \begin{math}[Y]_{4096}\end{math}，或一个十六进制的常数，比如 [0x0C]\begin{math}_{8}\end{math}。根据需要，\begin{math}[x]_{t}\end{math} 右边可以加上\begin{math} (s-t)\end{math} 个 0，使字串长度扩展成 \begin{math}s\end{math} 位，得到 \begin{math}[x]_{s}\end{math}。例如：
       [0x05]\begin{math}_{8} = 00000101\end{math}，
       [0x05]\begin{math}_{16} = 00000101 00000000\end{math}，
       [0x0005]\begin{math}_{16} = 00000000 00000101\end{math}，
       [0x13]\begin{math}_{8} = 00010011\end{math}，
       [0x13]\begin{math}_{16} = 00010011 00000000\end{math}，
       [0x0013]\begin{math}_{16} = 00000000 00010011\end{math}。

    \item {\bfseries ||}:
       表示位串粘接操作符，用于将两个位串前后粘成一个较长的位串。
\end{itemize}



\subsection{私钥运算子}
\label{sec:digsig-private-key-operands}

私钥运算子 \begin{math}Y\end{math}（私钥指数）和 \begin{math}M\end{math}（密钥模数）由用户生成。它们具有特定的 RSA 密钥长度（最大为 4096 位）。
RSA 签名操作还额外需要两个运算子，即参数 \begin{math}\overline{r}\end{math} 和 \begin{math}M^{\prime}\end{math}，由 \begin{math}Y\end{math} 和 \begin{math}M\end{math} 运算得到。


用户加密运算子 \begin{math}Y\end{math}、
\begin{math}M\end{math}、
\begin{math}\overline{r}\end{math}、
\begin{math}M^{\prime}\end{math} 与验证摘要，并以密文 \begin{math}C\end{math} 的形式存储。密文 \begin{math}C\end{math} 输入到 RSA\_DS 外设之后，先由硬件解密，然后参与 RSA 签名运算。有关生成 \begin{math}C\end{math} 的方法，请参考章节 \ref{sec:digsig-prerequisities}。

RSA\_DS 外设支持运算子长度为 \begin{math}N = 32 \times x\end{math} (\begin{math}x \in \{ 1,2,3, \dots, 128\} \end{math}) 的 RSA 签名运算 \begin{math}Z = X^{Y} \bmod M\end{math}。其中，参数的位宽需要满足 RSA 计算的运算子长度要求，即 \begin{math}Z\end{math}、
\begin{math}X\end{math}、
\begin{math}Y\end{math}、
\begin{math}M\end{math} 和
\begin{math}\overline{r}\end{math} 的位宽必须相同，且为 \begin{math}N\end{math} 中的任意值，
而 \begin{math}M^{\prime}\end{math} 的位宽则始终是 32。
更多 RSA 计算相关信息，请参考章节 \ref{mod:rsa} \textit{\nameref{mod:rsa}} 中的 \ref{sec:rsa-large-modular-Exponentiation} \textit{\nameref{sec:rsa-large-modular-Exponentiation}} 部分。


\subsection{软件准备工作}
\label{sec:digsig-prerequisities}

如果用户想使用 RSA\_DS 外设进行数字签名，则需要提前做好软件准备工作，如图 \ref{fig:digsig-flow} 所示。图中左半部分列出了在硬件开始 RSA 签名计算之前软件需要做的准备工作，右半部分则列出了硬件在整个签名计算过程中的计算流程。


\begin{figure}[H]
    \centering
    \includegraphics[width=0.9\textwidth]{\modulefiles/figures/31-RSADS-flow__CN}
    \caption{软件准备工作与硬件工作流程}
    \label{fig:digsig-flow}
\end{figure}


\begin{tiplisting}
\vspace{-2.5em}
\begin{enumerate}
\item 即便有多次签名计算，也只需进行一次软件准备工作（图 \ref{fig:digsig-flow} 左侧），但每一次签名计算都需要重新进行硬件运算（图 \ref{fig:digsig-flow} 右侧）。
\item 软件准备需要配置时钟复位，相关信息请参考章节 \ref{mod:resclk} \textit{\nameref{mod:resclk}}。
\end{enumerate}
\end{tiplisting}


用户需要依照图 \ref{fig:digsig-flow} 左侧中指定的步骤来计算 \begin{math}C\end{math}。详细过程如下：

\begin{itemize}

    \item {\bfseries 步骤 1}：按照章节 \ref{sec:digsig-private-key-operands} 所述，准备好符合运算子长度要求的 RSA 私钥运算子 \begin{math}Y\end{math} 和 \begin{math}M\end{math}。
          定义 \begin{math}[L]_{32} = \frac{N}{32}-1 \end{math}（比如，对于 RSA 4096，\begin{math}[L]_{32}\end{math} == [0x80-1]\begin{math}_{32}\end{math}）。根据私钥的生成方法准备 \begin{math}[RSA\_DS\_KEY]_{256}\end{math}：
          \begin{itemize}
              \item 当使用 HMAC 作为密钥来源时，准备好 \begin{math}[HMAC\_KEY]_{256}\end{math}，并计算出 \begin{math}[RSA\_DS\_KEY]_{256}\end{math}， 即 \begin{math}RSA\_DS\_KEY\end{math} = HMAC-SHA256(\begin{math}[HMAC\_KEY]_{256}\end{math}, \begin{math}1^{256}\end{math})。
              \item 当使用密钥管理器作为密钥来源时，准备好 256 位的 \begin{math}RSA\_DS\_KEY\end{math}。
          \end{itemize}
          接着，生成一个符合 AES-CBC 块加密算法要求的随机 \begin{math}[IV]_{128}\end{math}。有关 AES 的更多信息，请参考章节 \ref{mod:aes} \textit{\nameref{mod:aes}}。

    \item {\bfseries 步骤 2}：
          根据 \begin{math}M\end{math} 求解 \begin{math}\overline{r}\end{math} 和 \begin{math}M^{\prime} \end{math}。

    \item {\bfseries 步骤 3}：
          扩展 \begin{math}Y\end{math}、
               \begin{math}M\end{math} 和
               \begin{math}\overline{r}\end{math}，
          得到 \begin{math}[Y]_{4096}\end{math}、
               \begin{math}[M]_{4096}\end{math} 和
               \begin{math}[\overline{r}]_{4096}\end{math}。
          此步骤是由于 \begin{math}Y\end{math}、\begin{math}M\end{math} 和 \begin{math}\overline{r}\end{math} 的最大位宽为 4096，故位宽小于 4096 的运算子需要扩展至 4096，而位宽等于 4096 的则不需要扩展。


    \item {\bfseries 步骤 4}：
            使用 SHA-256 计算 MD 校验码：
                \begin{math}[MD]_{256}\end{math} = SHA256 (\begin{math}
                [Y]_{4096} ||
                [M]_{4096} ||
                [\overline{r}]_{4096} ||
                [M^{\prime}]_{32} ||
                [L]_{32} ||
                [IV]_{128}
            \end{math})

    \item {\bfseries 步骤 5}：构造 \begin{math}[P]_{12672}\end{math} = (
            \begin{math}
                [Y]_{4096} ||
                [M]_{4096} ||
                [\overline{r}]_{4096} ||
                [Box]_{384}
            \end{math})。其中
            \begin{math}[Box]_{384}\end{math} = (
            \begin{math}
                [MD]_{256} ||
                [M^{\prime}]_{32} ||
                [L]_{32} ||
                [\beta]_{64})
            \end{math}，
            \begin{math}[\beta]_{64} \end{math} 为 PKCS\#7 填充值，即由 8 个字节（0x08）组成的 $[$0x0808080808080808$]_{64}$ 字符串。\begin{math}[\beta]_{64}\end{math} 的作用是使 \begin{math}P\end{math} 的位长度为 128 的整数倍。
    \item {\bfseries 步骤 6}：
            计算密文形式的私钥参数 \begin{math}C\end{math} = \begin{math}[C]_{12672}\end{math} =
            AES-CBC-ENC (\begin{math}[P]_{12672}\end{math}, \begin{math}[RSA\_DS\_KEY]_{256}\end{math}, \begin{math}[IV]_{128}\end{math})，其中 $C$ 为密文，长度为 12672 位。
            \begin{math}C\end{math} 也可以表示为
            \begin{math}C\end{math} = \begin{math}[C]_{12672}\end{math} =
            (\begin{math}
            [\widehat{Y}]_{4096} ||
            [\widehat{M}]_{4096} ||
            [\widehat{ \overline{r} }]_{4096} ||
            [\widehat{Box}]_{384}
            \end{math})。
            其中，
            \begin{math} [\widehat{Y}]_{4096} \end{math}、
            \begin{math} [\widehat{M}]_{4096} \end{math}、
            \begin{math} [\widehat{ \overline{r} }]_{4096} \end{math}、
            \begin{math} [\widehat{Box}]_{384} \end{math}
            是 \begin{math}C\end{math} 的四个子参数，分别对应
            \begin{math} [Y]_{4096} \end{math}、
            \begin{math} [M]_{4096} \end{math}、
            \begin{math} [\overline{r}]_{4096} \end{math}、
            \begin{math} [Box]_{384} \end{math}
            的密文形式。

\end{itemize}

\subsection{硬件工作流程}
\label{sec:digsig-hardware-level-operation}

每次需要计算 RSA 数字签名时，都会触发硬件操作。硬件需要三个输入信息：预先生成的私钥密文 \begin{math}C\end{math}、唯一消息 \begin{math}X\end{math} 以及 \begin{math}IV\end{math}。

RSA\_DS 外设的工作流程可以分为三个阶段：

\begin{enumerate}

    \item {\bfseries 解析阶段，即图 \ref{fig:digsig-flow} 中的步骤 7 和 步骤 8}

        解析过程是图 \ref{fig:digsig-flow} 中步骤 6 的逆过程。RSA\_DS 外设将调用 AES 硬件加速器以 CBC 块模式对输入的密文信息 \begin{math}C\end{math} 进行解密，获取明文信息。
        该过程可以表示为
        \begin{math}P\end{math} = AES-CBC-DEC (\begin{math}C\end{math}, \begin{math}RSA\_DS\_KEY\end{math}, \begin{math}IV\end{math})。其中， \begin{math}IV\end{math} 就是 \begin{math}[IV]_{128}\end{math}，由用户直接指定； \begin{math}[RSA\_DS\_KEY]_{256}\end{math} 由：
        \begin{itemize}
            \item 硬件 HMAC 提供，通过存储在 eFuse 中的 \begin{math}HMAC\_KEY\end{math} 得到，软件无法获取。
            \item 硬件密钥管理器提供时，通过部署在密钥管理器中的 \begin{math}RSA\_DS\_KEY\end{math} 得到，软件无法获取。
        \end{itemize}
        通过 \begin{math}P\end{math}，RSA\_DS 外设能够解析出 \begin{math}[Y]_{4096}\end{math}、\begin{math}[M]_{4096}\end{math}、\begin{math}[\overline{r}]_{4096}\end{math}、\begin{math}[M^{\prime}]_{32}\end{math}、\begin{math}[L]_{32}\end{math}、MD 校验码和追加码 \begin{math}[\beta]_{64} \end{math}，这相当于步骤 5 的逆过程。

    \item {\bfseries 校验阶段，即图 \ref{fig:digsig-flow} 中的步骤 9 和步骤 10}

        RSA\_DS 外设会执行两种校验操作： MD 校验和填充 (padding) 校验。
        由于填充校验和 MD 校验同步进行，因此填充校验未在图 \ref{fig:digsig-flow} 中体现。


        \begin{itemize}

            \item MD 校验：RSA\_DS 外设调用 SHA-256 进行哈希计算获取哈希值 \begin{math}[CALC\_MD]_{256}\end{math}（计算方式与参数同 \begin{math}[MD]_{256}\end{math}，参见步骤 4）。然后，将 \begin{math}[CALC\_MD]_{256}\end{math} 与 MD 校验码 \begin{math}[MD]_{256}\end{math} 作比较，当且仅当二者相同时，MD 校验通过。

            \item 填充校验：RSA\_DS 外设检查解析阶段解析出的 追加码 \begin{math}[\beta]_{64} \end{math} 是否符合 PKCS\#7 标准，当且仅当符合标准时，填充校验通过。

        \end{itemize}

        仅当 MD 校验通过后，RSA\_DS 外设才会执行后续计算，否则 RSA\_DS 外设将拒绝执行。
        如果填充校验失败，将生成警告信息，但不会影响 RSA\_DS 外设的后续操作。

    \item {\bfseries 计算阶段，即图 \ref{fig:digsig-flow} 中的步骤 11 和步骤 12}

        RSA\_DS 外设将把用户输入的 \begin{math}X\end{math}，
        以及解析得到的 \begin{math}Y\end{math}、\begin{math}M\end{math} 和 \begin{math}\overline{r}\end{math} 都视为大数，结合解析得到的 \begin{math}M^{\prime}\end{math}，构成了大数模幂运算 \begin{math} X^{Y} \bmod M\end{math} 的所有必要输入参数。大数模幂运算的运算长度由 \begin{math}L\end{math} 的值唯一指定。RSA\_DS 外设调用 RSA 完成大数模幂运算 \begin{math}Z = X^{Y} \bmod M\end{math}，\begin{math}Z\end{math} 为签名结果。

\end{enumerate}

\subsubsection{硬件占用}

在运算过程中，\chipname{} 的 RSA\_DS 将占用以下硬件模块：
\begin{itemize}
    \item AES 加速器
    \item RSA 加速器
    \item SHA 加速器
\end{itemize}

每个加速器都有一个对应的 free 寄存器字段，用于指示 RSA\_DS 是否可以立即占用该加速器。如果 free 寄存器字段设置为 1，RSA\_DS 将在开始运算时立即复位该加速器。如果 free 寄存器字段设置为 0，RSA\_DS 将等待该加速器完成当前运算，或等到 \hyperref[fielddesc:RSADSTIMEOUTLIMIT]{RSA\_DS\_TIMEOUT\_LIMIT} 设置的超时上限。

当 RSA\_DS 开始运算并占用加速器后，各加速器的释放方式如下：
\begin{itemize}
    \item CBC 解密完成后，AES 加速器将被释放。
    \item RSA 加速器在整个 RSA\_DS 运算期间保持占用。
    \item MD 校验完成后，SHA 加速器将被释放。
\end{itemize}

如果 RSA\_DS 在占用一个或多个硬件加速器期间接收到外部复位，该复位也会同步传播至其占用的加速器。该机制确保暂存在这些加速器内部的 RSA\_DS 中间结果会被清除，从而防止其因外部复位攻击而发生泄漏。

\begin{tiplisting}
硬件占用是一种保护复用模块和存储空间的机制。当一个模块被硬件占用时，用户将无法执行以下操作：
\vspace{-0.5em}
\begin{itemize}
    \item 读取或写入模块的寄存器或存储器数据。
    \item 关闭模块时钟。
    \item 复位模块。
\end{itemize}
\vspace{-0.5em}
在硬件占用结束时，占用的模块将自动复位。此外，当用户对主模块执行软件复位时，所有占用的模块将同时复位。
\end{tiplisting}

\subsection{软件工作流程}

每次需要计算数字签名时，都应执行以下软件操作。输入内容为预先生成的私钥密文 \begin{math}C\end{math}、唯一的消息 \begin{math}X\end{math}、 \begin{math}IV\end{math}。本节中的软件步骤会触发章节 \ref{sec:digsig-hardware-level-operation} 中描述的硬件工作流程。\\
下述流程基于一个假设：软件已经调用了 HMAC 模块，HMAC 已经根据 \begin{math}HMAC\_KEY\end{math} 计算出了 \begin{math}RSA\_DS\_KEY\end{math}。

\begin{enumerate}

    \item {\bfseries 准备工作}：准备好 \begin{math}C\end{math}、\begin{math}X\end{math}、 \begin{math}IV\end{math}。具体方法请参考章节 \ref{sec:digsig-prerequisities} 。
    \begin{itemize}
        \item 若需禁用 PKCS\#7 填充校验，请将 \hyperref[regdesc:RSADSSTATICREG]{RSA\_DS\_STATIC\_REG} 中的 \hyperref[fielddesc:RSADSSKIPPADDINGCHECK]{RSA\_DS\_SKIP\_PADDING\_CHECK} 置为 1。
    \end{itemize}

    \item {\bfseries 启动 RSA\_DS}：对寄存器 \hyperref[regdesc:RSADSSETSTARTREG]{RSA\_DS\_SET\_START\_REG} 写 1。

    \item {\bfseries 选择 \begin{math}RSA\_DS\_KEY\end{math} 的来源}：配置 \hyperref[regdesc:RSADSSTATICREG]{RSA\_DS\_STATIC\_REG} 中的 \hyperref[fielddesc:RSADSKEYSOURCE]{RSA\_DS\_KEY\_SOURCE} 域，并检查是否已经准备好：
    \begin{itemize}
        \item 将 \hyperref[fielddesc:RSADSKEYSOURCE]{RSA\_DS\_KEY\_SOURCE} 置为 1，即选择密钥管理器中部署的 \begin{math}RSA\_DS\_KEY\end{math}。\\
        检查 \begin{math}RSA\_DS\_KEY\end{math} 已经成功部署，检查方式为：读取寄存器域 \hyperref[fielddesc:KEYMNGKEYDSVLD]{KEYMNG\_KEY\_DS\_VLD}，当其值为 1 时，代表 \begin{math}RSA\_DS\_KEY\end{math} 已经成功部署并准备好。
        \item 将 \hyperref[fielddesc:RSADSKEYSOURCE]{RSA\_DS\_KEY\_SOURCE} 置为 0，即选择从 HMAC 模块计算 \begin{math}RSA\_DS\_KEY\end{math}。\\
        轮询寄存器 \hyperref[regdesc:RSADSQUERYBUSYREG]{RSA\_DS\_QUERY\_BUSY\_REG} 直到读到 0。

        如果 \hyperref[regdesc:RSADSQUERYBUSYREG]{RSA\_DS\_QUERY\_BUSY\_REG} 超过 1 ms 还没读到 0，则说明 HMAC 初始化出现了问题。此时，软件应当读寄存器 \hyperref[regdesc:RSADSQUERYKEYWRONGREG]{RSA\_DS\_QUERY\_KEY\_WRONG\_REG}，根据返回值判断具体是下述哪一种情况：
        \begin{itemize}

            \item 如果读到零值，说明 HMAC 未被调用。

            \item 如果读到非零值 (1 \verb+~+ 15)，则说明 HMAC 被调用过，但是 RSA\_DS 外设没有成功从 HMAC 模块获取到 \begin{math}RSA\_DS\_KEY\end{math}。此种情况可能是因为受到了其他程序的干扰。

        \end{itemize}


    \end{itemize}

    \item {\bfseries 将 \begin{math}IV\end{math} 写入存储器 \hyperref[RSADSIVMEM]{RSA\_DS\_IV\_MEM}}：将 \begin{math}IV\end{math} block 中的内容写入存储器 \hyperref[RSADSIVMEM]{RSA\_DS\_IV\_MEM}（大小为 16 字节）。有关 \begin{math}IV\end{math} block 的更多信息，请参考章节 \ref{mod:aes} \textit{\nameref{mod:aes}}。

    \item {\bfseries 将 \begin{math}X\end{math} 写入存储器 \hyperref[RSADSXMEM]{RSA\_DS\_X\_MEM}}：将 \begin{math}X_{i}\end{math}
        (\begin{math}i \in \{ 0, 1, \dots, n-1 \} \end{math})
        写入存储器 \hyperref[RSADSXMEM]{RSA\_DS\_X\_MEM}，容量为 128 个字 (word)，其中 \begin{math} n = \frac{N}{32} \end{math} 。
        每一个字刚好存放一个 \begin{math}b\end{math} 进制数。存储块以小端字节序进行存储，即存储器的低地址存放运算子的低位进制数，高地址存放运算子的高位进制数。当 \begin{math}X\end{math} 的长度小于 128 个字时，存储器 \hyperref[RSADSXMEM]{RSA\_DS\_X\_MEM} 中有一部分空间未使用，该部分空间中的数据可以是任意值。
    \item {\bfseries 将 \begin{math}C\end{math} 写入存储器}：即将 \begin{math}C\end{math} 中的四个子参数分别写入对应的存储器。
    \begin{itemize}

        \item 将 \begin{math}\widehat{Y}_{i}\end{math}
        (\begin{math}i \in \{ 0, 1, \dots, 127 \} \end{math})
        写入存储器 \hyperref[RSADSYMEM]{RSA\_DS\_Y\_MEM}。

        \item 将 \begin{math}\widehat{M}_{i}\end{math}
        (\begin{math}i \in \{ 0, 1, \dots, 127 \} \end{math})
        写入存储器 \hyperref[RSADSMMEM]{RSA\_DS\_M\_MEM}。

        \item 将 \begin{math}\widehat{\overline{r}}_{i}\end{math}
        (\begin{math}i \in \{ 0, 1, \dots, 127 \} \end{math})
        写入存储器 \hyperref[RSADSRBMEM]{RSA\_DS\_RB\_MEM}。

        \item 将 \begin{math}\widehat{Box}_{i}\end{math}
        (\begin{math}i \in \{ 0, 1, \dots, 11 \} \end{math})
        写入存储器 \hyperref[RSADSBOXMEM]{RSA\_DS\_BOX\_MEM}。

        其中 存储器
        \hyperref[RSADSYMEM]{RSA\_DS\_Y\_MEM}、
        \hyperref[RSADSMMEM]{RSA\_DS\_M\_MEM}、
        \hyperref[RSADSRBMEM]{RSA\_DS\_RB\_MEM} 的
        容量均为 128 个字，
        而存储器 \hyperref[RSADSBOXMEM]{RSA\_DS\_BOX\_MEM} 容量只有 12 个字。
        每一个字刚好存放一个 \begin{math}b\end{math} 进制数。
        存储块以小端字节序进行存储，即存储器的低地址存放运算子的低位进制数，高地址存放运算子的高位进制数。

    \end{itemize}

    \item {\bfseries 启动计算}：对寄存器 \hyperref[regdesc:RSADSSETCONTINUEREG]{RSA\_DS\_SET\_CONTINUE\_REG} 写入 1。
    \item {\bfseries 等待运算结束}：轮询寄存器 \hyperref[regdesc:RSADSQUERYBUSYREG]{RSA\_DS\_QUERY\_BUSY\_REG} 直到读到 0。

    \item {\bfseries 检查校验结果}：读寄存器 \hyperref[regdesc:RSADSQUERYCHECKREG]{RSA\_DS\_QUERY\_CHECK\_REG}，根据返回值决定后续操作。

    \begin{itemize}

        \item 如果返回值为 0，则说明填充校验通过，MD 校验通过，可以继续读取 \begin{math}Z\end{math} 结果值。

        \item 如果返回值为 1，则说明填充校验通过，但 MD 校验失败。\begin{math}Z\end{math} 结果值无效，跳至步骤 \ref{other:digsig-step-exit}。

        \item 如果返回值为 2，则说明填充校验失败，但 MD 校验通过，用户可以继续读取 \begin{math}Z\end{math} 结果值。但仍需注意的是，此时数据填充不符合 PKCS\#7 封装方式。

        \item 如果返回值为 3，则说明填充校验失败，且 MD 校验失败。 这种情况说明有致命错误发生。\begin{math}Z\end{math} 结果值无效。跳至步骤 \ref{other:digsig-step-exit}。

    \end{itemize}

    \item {\bfseries 读出运算结果}：从存储器 \hyperref[RSADSZMEM]{RSA\_DS\_Z\_MEM} 读出运算结果 \begin{math}Z_{i}\end{math} (\begin{math}i \in \{0,1, \dots,n-1\}\end{math})，其中 \begin{math} n = \frac{N}{32} \end{math} 。\begin{math}Z\end{math} 以小端字节序存储在存储器中。

    \item {\bfseries 退出计算环境\label{other:digsig-step-exit}}：对寄存器 \hyperref[regdesc:RSADSSETFINISHREG]{RSA\_DS\_SET\_FINISH\_REG} 写入 1，然后轮询寄存器 \hyperref[regdesc:RSADSQUERYBUSYREG]{RSA\_DS\_QUERY\_BUSY\_REG} 直到读到 0。

\end{enumerate}

运算结束后，所有输入/输出寄存器和存储器中的数据均已被清除。


%%%%%%%%%%%%%%%%%%%%%%%%
%%%%% Memory Block %%%%%
%%%%%%%%%%%%%%%%%%%%%%%%

\newpage
\section{ 存储器列表 }
\label{sec:digsig-mem-summ}

请注意，这里的起始地址和结束地址都是相对于 RSA 数字签名外设 (RSA\_DS) 基地址的地址偏移量（相对地址），具体基地址请见章节 \ref{mod:sysmem} \textit{\nameref{mod:sysmem}} 中的表 \ref{tab:sysmem-base-address}。

\subfile{\modulefiles/31-RSADS-mem-summ__CN}


%%%%%%%%%%%%%%%%%%%%%%%%
%%% Register Summary %%%
%%%%%%%%%%%%%%%%%%%%%%%%

% >>> Update the sample text below and insert
%     the info on register summary into the subfile <<<

\newpage
\hypertarget{digsig-reg-summ}{}
\section{寄存器列表}
\label{sec:digsig-reg-summ}

本小节的所有地址均为相对于 RSA 数字签名外设 (RSA\_DS) 基地址的地址偏移量（相对地址），具体基地址请见章节 \ref{mod:sysmem} \textit{\nameref{mod:sysmem}} 中的表 \ref{tab:sysmem-base-address}。

请查看章节 \hyperref[glossary-access-types]{\textit{寄存器的访问类型}}，了解“访问”列缩写的含义。

\subfile{\modulefiles/31-RSADS-reg-summ__CN}


%%%%%%%%%%%%%%%%%
%%% Registers %%%
%%%%%%%%%%%%%%%%%

% >>> Update the sample text below and insert
%      the info on registers into the subfile <<<

\newpage
\section{寄存器}

本小节的所有地址均为相对于 RSA 数字签名外设 (RSA\_DS) 基地址的地址偏移量（相对地址），具体基地址请见章节 \ref{mod:sysmem} \textit{\nameref{mod:sysmem}} 中的表 \ref{tab:sysmem-base-address}。

\subfile{\modulefiles/31-RSADS-reg__CN}


\end{document}
